\section{Conclusion: Closure of the PR Equation Architecture}
\label{sec:conclusion}

PR-IV establishes the requested closing statement for the admissible PR-I through PR-III equation architecture.  The theorem domain is fixed before any diagnostic comparison: the public finite-tier ledger is inherited, and every equation-selection coordinate must satisfy inheritance, covariance,
minimality, completeness, nonredundancy, canonical normalization, and the no-fit firewall.  Gauge, normalized-basis, coordinate, phase, algebraic, field-redefinition, and unit changes are quotiented as equivalent presentations.

The exact normal forms have sole admissible values.  The radial coefficients are $(a_2,a_4)=(1,3/4)$; the minimal required-channel projector retains $\operatorname{span}\{\phi_1,\phi_3\}$; signed terminal incidence fixes $\sigma_8=-1$; the nonzero hypercharge ledger leaves one orbit; the fixed
neutral mass operator leaves one photon null line; normalized complete and nonredundant sheet incidence fixes $(N_{\rm gen},[\Phi_3])=(3,[I_3])$; and the neutral response fixes $k=2$.
For the last coordinate, the structural tensor is
\begin{equation}
 D_3(k)=\rho kP_Z,
\end{equation}
while unitary-basis invariance, orthogonal-mode additivity, quadratic homogeneity, and unit-charge normalization uniquely force
\begin{equation}
 \mathcal F(q)=\operatorname{Tr}(q^\dagger q).
\end{equation}
The normalized charged-adjoint spectrum is $(+1,-1)$, so
$\mathcal F(q)=2$ and
\begin{equation}
 \boxed{D_3=2\rho P_Z.}
\end{equation}
This is not a best-fit choice: every $k\ne2$ has the exact positive defect
$\rho^2(k-2)^2$.

Combining the seven unique fibers in dependency order with the fixed inherited PR-I through PR-III maps gives
\begin{equation}
 \boxed{
 \mathfrak A_{\rm PR}^{\rm can}/\!\sim
 =\{[\mathrm{PR}]\}.}
\end{equation}
There is therefore no second inequivalent equation, branch, normalization, or continuous parameter allowed by the complete PR admissibility guidelines.  A differently written formula that works on the complete declared domain is a ledger-preserving presentation of the same physical equation.  An inequivalent in-class proposal has a nonzero exact defect.  A proposal that adds an operator, field, scale, Wilson \cite{Wilson_1975} coefficient, response jet, higher order, or different boundary/representation rule violates the class definition and is a different theory rather than another PR solution.

The second-jet and effective-field-theory constructions displayed in the appendices make that last distinction explicit.  They produce freedom only by dropping canonical normalization or by adding independent operator data. They therefore have a named nonzero admissibility defect and cannot enlarge
the PR quotient.  A determinant expression satisfying the same normalized response conditions is an equivalent microscopic representation of the already unique $D_3$ operator, not an additional coefficient or branch.

Direct postulate derivation and canonical logical closure are the two accepted provenance routes.  Neither route imports empirical targets, and neither adds an adjustable postulate.  The completed class has no unresolved selector and requires no residual ranking to choose among equations.  PR-IV therefore closes the mathematical question posed here: under the developed PR guidelines, all admissible formulations belong to one equation-equivalence class and no other PR solution architecture is allowed.
